\documentclass[border=5mm]{standalone}
% ==============================================================================
% SETUP.TEX - CONFIGURACIÓN GLOBAL DEL PROYECTO
% ==============================================================================
% Este archivo centraliza todos los paquetes y estilos.
% Debe incluirse en el preámbulo del libro principal y de los archivos standalone.
% \PassOptionsToPackage{gray}{xcolor}
% --- 1. CODIFICACIÓN E IDIOMA ---
% fix-cm habilita las fuentes Computer Modern en tamaños arbitrarios (por debajo
% de 5 pt, entre otros). Debe cargarse antes que fontenc.
\usepackage{fix-cm}
\usepackage[utf8]{inputenc}
\usepackage[T1]{fontenc}
\usepackage[spanish, es-tabla]{babel}  
\usepackage{csquotes} % Recomendado para babel y citas

\usepackage{import}
% mode=image: Usa el PDF pre-compilado, nunca intenta recompilar.
% Debes compilar cada figura manualmente antes de compilar el libro.
\usepackage[mode=image, subpreambles=false]{standalone}

% --- 2. MATEMÁTICAS Y CIENCIAS ---
\usepackage{amsmath, amssymb, amsfonts, mathtools}
\usepackage{enumitem}

% LaTeX trae \DeclareMathSizes{5}{5}{5}{5}, así que a 5 pt (\tiny) los subíndices
% salen del mismo tamaño que la base: en las etiquetas de figuras el M de
% $\omega_M$ queda desproporcionado. Se restablece la proporción habitual.
\DeclareMathSizes{5}{5}{3.7}{3}

% --- 3. DISEÑO Y GEOMETRÍA --- 
% Unificamos los márgenes para todo el documento
\usepackage[left=2.5cm,right=2.5cm,top=3cm,bottom=3cm]{geometry}
\makeatletter
\ifstandalone % Evita que geometry dañe el recorte de las figuras bajándolas 3cm
  \@ifclasswith{standalone}{crop=false}{
    % Es un capítulo/sección: Dejamos que geometry actúe.
  }{
    % Es una figura: Neutralizamos geometry para que no las recorte mal.
    \geometry{pass}
  }
\fi
\makeatother
\usepackage{parskip} % Espaciado entre párrafos sin sangría
\usepackage{float}   % Para usar [H] en figura

% Ajustes de flotantes para evitar espacios verticales exagerados
\renewcommand{\topfraction}{0.95}      % Permite que las figuras ocupen hasta el 95% superior
\renewcommand{\bottomfraction}{0.95}   % Permite que las figuras ocupen hasta el 95% inferior
\renewcommand{\textfraction}{0.05}     % Permite hojas con tan solo un 5% de texto
\renewcommand{\floatpagefraction}{0.8} % Requiere que una página de solo figuras esté 80% llena
\setcounter{topnumber}{4}
\setcounter{bottomnumber}{4}
\setcounter{totalnumber}{10}

% --- 4. GRÁFICOS Y TABLAS ---
\usepackage{graphicx}
\usepackage{booktabs} % Tablas profesionales
\usepackage{caption}  % Control de captions y \captionof
\usepackage{tikz}
\usetikzlibrary{arrows.meta, calc, angles, quotes, babel, backgrounds}
\usepackage{pgfplots}
\usepgfplotslibrary{groupplots}
\usepgfplotslibrary{external}
\usepgfplotslibrary{patchplots}
\pgfplotsset{compat=1.18}
\usepackage[american, straightvoltages]{circuitikz}

% --- 5. COLORES Y CAJAS ---

\usepackage[table]{xcolor}
\usepackage{tcolorbox}

% Definición de colores corporativos/estilo
\definecolor{utnblue}{RGB}{0, 51, 102}
\definecolor{codegreen}{rgb}{0,0.6,0}
\definecolor{codegray}{rgb}{0.5,0.5,0.5}
\definecolor{codepurple}{rgb}{0.58,0,0.82}
\definecolor{backcolour}{rgb}{0.96,0.96,0.96}

% --- 6. ENTORNOS PERSONALIZADOS (CAJAS) ---

% Caja "Resultado" (Azul Corporativo) — Definiciones, fórmulas clave, resultados importantes.
% Uso: \begin{resultado}[Fórmula de Euler] ... \end{resultado}
\newtcolorbox{resultado}[1][]{
    colback=utnblue!5!white,
    colframe=utnblue,
    title=\textbf{#1},
    fonttitle=\bfseries,
    sharp corners,
    boxrule=0.5mm
}

% Caja "Nota" (Gris) — Advertencias, aplicaciones, contexto secundario.
% Uso: \begin{nota}[Cuidado con la calculadora] ... \end{nota}
% Uso: \begin{nota} ... \end{nota}  → título por defecto "Nota"
\newtcolorbox{nota}[1][Nota]{
    colback=codegray!5!white,
    colframe=codegray,
    title=\textbf{#1},
    fonttitle=\bfseries,
    sharp corners,
    boxrule=0.5mm
}

% Caja inline para resaltar un valor y enlazarlo con su otra aparición en una tabla
% o en una cuenta: mismo color, mismo valor.
% Uso: \marcavalor{$4$}  o  \marcavalor[codegreen]{$4$}
\newtcbox{\marcavalor}[1][utnblue]{
    on line,
    boxsep=1pt, left=2pt, right=2pt, top=1pt, bottom=1pt,
    colback=#1!10!white,
    colframe=#1,
    boxrule=0.4pt,
    arc=2pt
}

% --- 7. CÓDIGO FUENTE (LISTINGS) ---
\usepackage{listings}

\lstdefinestyle{miestilo}{
    backgroundcolor=\color{backcolour},   
    commentstyle=\color{codegreen},
    keywordstyle=\color{magenta},
    numberstyle=\tiny\color{codegray},
    stringstyle=\color{codepurple},
    basicstyle=\ttfamily\footnotesize,
    breakatwhitespace=false,         
    breaklines=true,                 
    captionpos=b,                    
    keepspaces=true,                 
    numbers=left,                    
    numbersep=5pt,                  
    showspaces=false,                
    showstringspaces=false,
    showtabs=false,                  
    tabsize=2,
    frame=single,                    
    rulecolor=\color{codegray}
}

\lstset{style=miestilo}
% Soporte para tildes y ñ en bloques de código
\lstset{literate={ñ}{{\~n}}1 {á}{{\'a}}1 {é}{{\'e}}1 {í}{{\'i}}1 {ó}{{\'o}}1 {ú}{{\'u}}1}

% Operadores matemáticos personalizados

% Vector en sentido discreto (lista finita de coordenadas): negrita + flecha.
% Uso: \vect{v}, \vect{e}_k, \vect{0}.
\newcommand{\vect}[1]{\vec{\mathbf{#1}}}

% Fecha de versión y comando para capítulos standalone
\providecommand{\verdate}{\the\year.\ifnum\month<10 0\fi\the\month.\ifnum\day<10 0\fi\the\day}
\newcommand{\chapterversion}{%
    \vspace{-1.5cm}\begin{flushright}\small\textit{\color{codegray}Versión~\verdate}\end{flushright}\vspace{0.5cm}%
}

% --- 8. HIPERVÍNCULOS (DEBE SER EL ÚLTIMO PAQUETE) ---
\usepackage[colorlinks=true, linkcolor=utnblue, urlcolor=utnblue, citecolor=utnblue]{hyperref}

% --- 9. REFERENCIAS CRUZADAS ENTRE CAPÍTULOS STANDALONE ---
% Cuando se compila un capítulo standalone, xr-hyper lee los .aux de los
% otros capítulos (si existen) para resolver \ref{} a labels externos.
% En la compilación del libro completo este bloque no se activa.
\ifstandalone
  \usepackage{xr-hyper}
  \makeatletter
  \newcommand{\xrchap}[1]{%
    \edef\xrc@self{\detokenize{Capitulo#1}}%
    \edef\xrc@job{\jobname}%
    \ifx\xrc@self\xrc@job\else
      \IfFileExists{../#1capitulo/Capitulo#1.aux}%
        {\externaldocument{../#1capitulo/Capitulo#1}}{}%
    \fi
  }
  \makeatother
  \xrchap{01}\xrchap{02}\xrchap{03}\xrchap{04}
  \xrchap{05}\xrchap{06}\xrchap{07}\xrchap{08}
  \xrchap{09}\xrchap{10}\xrchap{11}\xrchap{12}
  \xrchap{13}
\fi
% \PassOptionsToPackage{gray}{xcolor}

\begin{document}
    \begin{tikzpicture}[>=Latex, font=\small]
        
        % Señal base: sin(1.5*x) + 0.3*sin(4*x)
        % Todos los valores calculados con Python para coherencia

        \pgfplotsset{
            width=7cm, height=4cm,
            axis lines=center,
            axis line style={->, thick},
            xlabel style={anchor=north west, at={(axis description cs:1,0)}, font=\footnotesize},
            ylabel style={anchor=south east, at={(axis description cs:0,1)}, font=\footnotesize},
            xmin=0, xmax=6, ymin=-1.5, ymax=1.5,
            ticks=none,
        }

        % =============================================
        % (a) Continua en t, Continua en amplitud (Analógica)
        % =============================================
        \begin{axis}[
            at={(0cm,0cm)}, anchor=center,
            xlabel={$t$}, ylabel={$x(t)$},
        ]
            \addplot[domain=0:6, samples=300, ultra thick, utnblue, smooth] 
                {sin(deg(1.5*x)) + 0.3*sin(deg(4*x))};
        \end{axis}
        \node[font=\footnotesize\bfseries] at (0,2.5) {(a) Analógica};
        \node[font=\tiny, gray] at (0,2) {Tiempo continuo, Amplitud continua};

        % =============================================
        % (b) Continua en t, Cuantizada en amplitud
        % =============================================
        \begin{axis}[
            at={(8cm,0cm)}, anchor=center,
            xlabel={$t$}, ylabel={$x_q(t)$},
        ]
            % Niveles cuantización
            \draw[gray!20, thin] (axis cs:0,-1.25) -- (axis cs:6,-1.25);
            \draw[gray!20, thin] (axis cs:0,-1.0) -- (axis cs:6,-1.0);
            \draw[gray!20, thin] (axis cs:0,-0.75) -- (axis cs:6,-0.75);
            \draw[gray!20, thin] (axis cs:0,-0.5) -- (axis cs:6,-0.5);
            \draw[gray!20, thin] (axis cs:0,-0.25) -- (axis cs:6,-0.25);
            \draw[gray!20, thin] (axis cs:0,0) -- (axis cs:6,0);
            \draw[gray!20, thin] (axis cs:0,0.25) -- (axis cs:6,0.25);
            \draw[gray!20, thin] (axis cs:0,0.5) -- (axis cs:6,0.5);
            \draw[gray!20, thin] (axis cs:0,0.75) -- (axis cs:6,0.75);
            \draw[gray!20, thin] (axis cs:0,1.0) -- (axis cs:6,1.0);
            \draw[gray!20, thin] (axis cs:0,1.25) -- (axis cs:6,1.25);
            % Curva escalonada — valores correctos de round(f(x)*4)/4
            \addplot[thick, red!70!black, const plot] coordinates {
                (0.00, +0.00) (0.15, +0.50) (0.30, +0.75) (0.45, +1.00) (0.60, +1.00)
                (0.75, +1.00) (0.90, +0.75) (1.05, +0.75) (1.20, +0.75) (1.35, +0.75)
                (1.50, +0.75) (1.65, +0.75) (1.80, +0.75) (1.95, +0.50) (2.10, +0.25)
                (2.25, +0.00) (2.40, -0.50) (2.55, -0.75) (2.70, -1.00) (2.85, -1.25)
                (3.00, -1.25) (3.15, -1.00) (3.30, -0.75) (3.45, -0.50) (3.60, -0.50)
                (3.75, -0.50) (3.90, -0.50) (4.05, -0.25) (4.20, -0.25) (4.35, +0.00)
                (4.50, +0.25) (4.65, +0.50) (4.80, +1.00) (4.95, +1.25) (5.10, +1.25)
                (5.25, +1.25) (5.40, +1.00) (5.55, +0.75) (5.70, +0.50) (5.85, +0.25)
                (6.00, +0.25)
            };
        \end{axis}
        \node[font=\footnotesize\bfseries] at (8,2.5) {(b) Cuantizada};
        \node[font=\tiny, gray] at (8,2) {Tiempo continuo, Amplitud cuantizada};

        % =============================================
        % (c) Discreta en t, Continua en amplitud (Muestreada)
        % =============================================
        \begin{axis}[
            at={(0cm,-5.5cm)}, anchor=center,
            xlabel={$n$}, ylabel={$x[n]$},
        ]
            % Señal subyacente gris
            \addplot[domain=0:6, samples=300, thin, gray!30, smooth] 
                {sin(deg(1.5*x)) + 0.3*sin(deg(4*x))};
            % Stems — valores exactos de f(x) = sin(1.5x) + 0.3*sin(4x)
            \addplot+[ycomb, mark=*, mark size=2pt, mark options={fill=utnblue, draw=utnblue},
                     thick, utnblue] coordinates {
                (0.0,  +0.0000) (0.3,  +0.7146) (0.6,  +0.9860) (0.9,  +0.8430)
                (1.2,  +0.6750) (1.5,  +0.6942) (1.8,  +0.6655) (2.1,  +0.2480)
                (2.4,  -0.4948) (2.7,  -1.0828) (3.0,  -1.1385) (3.3,  -0.7943)
                (3.6,  -0.4831) (3.9,  -0.3874) (4.2,  -0.2495) (4.5,  +0.2247)
                (4.8,  +0.8967) (5.1,  +1.2792) (5.4,  +1.0843) (5.7,  +0.5504)
            };
        \end{axis}
        \node[font=\footnotesize\bfseries] at (0,-3) {(c) Muestreada};
        \node[font=\tiny, gray] at (0,-3.5) {Tiempo discreto, Amplitud continua};

        % =============================================
        % (d) Discreta en t, Cuantizada en amplitud (Digital)
        % =============================================
        \begin{axis}[
            at={(8cm,-5.5cm)}, anchor=center,
            xlabel={$n$}, ylabel={$x_q[n]$},
        ]
            % Niveles cuantización
            \draw[gray!20, thin] (axis cs:0,-1.25) -- (axis cs:6,-1.25);
            \draw[gray!20, thin] (axis cs:0,-1.0) -- (axis cs:6,-1.0);
            \draw[gray!20, thin] (axis cs:0,-0.75) -- (axis cs:6,-0.75);
            \draw[gray!20, thin] (axis cs:0,-0.5) -- (axis cs:6,-0.5);
            \draw[gray!20, thin] (axis cs:0,-0.25) -- (axis cs:6,-0.25);
            \draw[gray!20, thin] (axis cs:0,0) -- (axis cs:6,0);
            \draw[gray!20, thin] (axis cs:0,0.25) -- (axis cs:6,0.25);
            \draw[gray!20, thin] (axis cs:0,0.5) -- (axis cs:6,0.5);
            \draw[gray!20, thin] (axis cs:0,0.75) -- (axis cs:6,0.75);
            \draw[gray!20, thin] (axis cs:0,1.0) -- (axis cs:6,1.0);
            \draw[gray!20, thin] (axis cs:0,1.25) -- (axis cs:6,1.25);
            % Stems cuantizados — round(f(x)*4)/4
            \addplot+[ycomb, mark=*, mark size=2pt, mark options={fill=red!70!black, draw=red!70!black},
                     thick, red!70!black] coordinates {
                (0.0,  +0.00) (0.3,  +0.75) (0.6,  +1.00) (0.9,  +0.75)
                (1.2,  +0.75) (1.5,  +0.75) (1.8,  +0.75) (2.1,  +0.25)
                (2.4,  -0.50) (2.7,  -1.00) (3.0,  -1.25) (3.3,  -0.75)
                (3.6,  -0.50) (3.9,  -0.50) (4.2,  -0.25) (4.5,  +0.25)
                (4.8,  +1.00) (5.1,  +1.25) (5.4,  +1.00) (5.7,  +0.50)
            };
        \end{axis}
        \node[font=\footnotesize\bfseries] at (8,-3) {(d) Digital};
        \node[font=\tiny, gray] at (8,-3.5) {Tiempo discreto, Amplitud cuantizada};

        % --- Flechas de operación ---
        % Horizontal superior: Cuantizar amplitud
        \draw[->, very thick, gray!60] (3.5,0.5) -- (4.5,0.5) 
            node[midway, above, font=\tiny\itshape] {Cuantizar};

        % Vertical izquierdo: Muestrear
        \draw[->, very thick, gray!60] (-3,-1.5) -- (-3,-2.5) 
            node[midway, left, font=\tiny\itshape, align=right] {Muestrear};

        % Vertical derecho: Muestrear
        \draw[->, very thick, gray!60] (11,-1.5) -- (11,-2.5) 
            node[midway, right, font=\tiny\itshape, align=left] {Muestrear};

        % Horizontal inferior: Cuantizar
        \draw[->, very thick, gray!60] (3.5,-5) -- (4.5,-5) 
            node[midway, above, font=\tiny\itshape] {Cuantizar};

        % Diagonal: Conversión A/D completa
        \draw[->, very thick, orange!70!black] (3.5,-0.5) -- (4.5,-4.5) 
            node[midway, right, font=\tiny\itshape, align=left, xshift=2mm] {Conversión\\A/D};

    \end{tikzpicture}
\end{document}
